\subsection{分配律的推论}

乘法对加法的分配律使我们能够应用 §3 第 4 号的命题 1，它给出

\[
\prod_ {i = 1} ^ {n} \left(\sum_ {\lambda \in L _ {i}} x _ {i, \lambda}\right) = \sum_ {\alpha_ {1}, \dots , \alpha_ {n}} \prod_ {i = 1} ^ {n} x _ {i, \alpha_ {i}}\tag{13}
\]

其中和式遍及所有属于 \((\alpha_{1},\ldots,\alpha_{n})\) 的序列 \(L_{1}\times\cdots\times L_{n}\) ，并且对于 \(i=1,\ldots,n\) ，环 A 的元素族 \((x_{i,\lambda})_{\lambda\in L^{i}}\) 具有有限支撑。

命题 1. 设 A 为交换环， \((x_{\lambda})_{\lambda \in L}\) 为 A 中元素的有限族。对于每个正整数族 \(\beta = (\beta_{\lambda})_{\lambda \in L}\) ，设 \(|\beta| = \sum_{\lambda \in L} \beta_{\lambda}\) 。则

\[
\left(\sum_ {\lambda \in L} x _ {\lambda}\right) ^ {n} = \sum_ {| \beta | = n} \frac {n !}{\prod_ {\lambda \in L} \beta_ {\lambda} !} \prod_ {\lambda \in L} x _ {\lambda} ^ {\beta_ {\lambda}}.\tag{14}
\]

我们应用公式 (13)，其中 \(\mathbf{L}_i = \mathbf{L}\) 和 \(x_{i,\lambda} = x_{\lambda}\) 对于 \(1 \leqslant i \leqslant n\) 。则

\[
\left(\sum_ {\lambda \in L} x _ {\lambda}\right) ^ {n} = \sum_ {\alpha_ {1}, \dots , \alpha_ {n}} x _ {\alpha_ {1}} \dots x _ {\alpha_ {n}},\tag{15}
\]

和式遍及所有序列 \(\alpha = (\alpha_{1},\ldots ,\alpha_{n})\in \mathbf{L}^{n}\) 。

设 \(\alpha\) 属于 \(\mathbf{L}^n\) ；对所有 \(\lambda \in \mathbf{L}\) ，设 \(\mathrm{U}_{\lambda}^{\alpha}\) 表示整数 \(i\) 组成的集合，满足 \(1 \leqslant i \leqslant n\) 且 \(\alpha_i = \lambda\) ，并设 \(\Phi(\alpha) = (\mathrm{U}_{\lambda}^{\alpha})_{\lambda \in \mathbf{L}}\) 。立即可见 \(\Phi\) 是 \(\mathbf{L}^n\) 到 \(\{1, 2, \ldots, n\}\) 的以 \(\mathbf{L}\) 为指标集的分拆组成的集合上的一个双射。对所有 \(\beta \in \mathbf{N}^{\mathbf{L}}\) （满足 \(|\beta| = n\) ），设 \(\mathbf{L}_{\beta}^n\) 表示由 \(\alpha \in \mathbf{L}^n\) （满足 \(\operatorname{Card} \mathrm{U}_{\beta}^{\alpha} = \beta_{\lambda}\) 对所有 \(\lambda \in \mathbf{L}\) 成立）组成的集合。由此可知族 \((\mathbf{L}_{\beta}^n)_{\beta | = n}\) 是 \(\mathbf{L}^n\) 的一个分拆，并且

\[
\operatorname{Card} \mathrm{L} _ {\beta} ^ {n} = \frac {n !}{\prod_ {\lambda \in \mathrm{L}} \beta_ {\lambda} !}
\]

（§ 5，第 5 号）。

最后，对 \(\alpha \in \mathbf{L}_{\beta}^{n}\) ，

\[
x _ {\alpha_ {1}} \dots x _ {\alpha_ {n}} = \prod_ {\lambda \in L} \prod_ {i \in U _ {\lambda} ^ {\alpha}} x _ {\alpha_ {i}} = \prod_ {\lambda \in L} \prod_ {i \in U _ {\lambda} ^ {\alpha}} x _ {\lambda} = \prod_ {\lambda \in L} x _ {\lambda},
\]

从而

\[
\begin{array}{r l} \sum_ {(\alpha_ {1}, \dots , \alpha_ {n})} x _ {\alpha_ {1}} \dots x _ {\alpha_ {n}} & = \sum_ {| \beta | = n} \sum_ {\alpha \in L _ {\beta} ^ {n}} x _ {\alpha_ {1}} \dots x _ {\alpha_ {n}} \\ & = \sum_ {| \beta | = n} \sum_ {\alpha \in L _ {\beta} ^ {n}} \prod_ {\lambda \in L _ {\lambda}} x _ {\lambda} ^ {\beta_ {\lambda}} \\ & = \sum_ {| \beta | = n} \frac {n !}{\prod_ {\lambda \in L} \beta_ {\lambda} !} x _ {\lambda} ^ {\beta_ {\lambda}} \end{array}
\]

从而公式 (14) 由 (15) 得出。

推论 1（二项式公式）。设 \(x\) 和 \(y\) 为交换环 A 的两个元素。则：

\[
(x + y) ^ {n} = \sum_ {p = 0} ^ {n} \binom {n} {p} x ^ {p} y ^ {n - p}.
\]

将公式 (14) 应用于 \(\mathbf{L} = \{1,2\}\) , \(x_{1} = x\) 和 \(x_{2} = y\) 给出

\[
(x + y) ^ {n} = \sum_ {p + q = n} \frac {n !}{p ! q !} x ^ {p} y ^ {q},
\]

其中和式遍及正整数有序对 \(p, q\) （满足 \(p + q = n\) ）。二项式公式由此直接得出（集合论，III，§ 5，第 8 号）。

推论 2. 设 A 为交换环，X 为一个集合， \(\mathbf{u} = (u_x)_{x \in \mathbb{X}}\) 和 \(\mathbf{v} = (v_x)_{x \in \mathbb{X}}\) 为 A 中元素的两个族。设 \(\mathbf{u} + \mathbf{v}\) 表示族 \((u_x + v_x)_{x \in \mathbb{X}}\) 。对所有 \(\lambda \in \mathbf{N}^{(\mathbb{X})}\) ，我们记 \(\lambda! = \prod_{x \in \mathbb{X}} \lambda(x)!\) 。则对所有 \(\alpha \in \mathbf{N}^{(\mathbb{X})}\) ，在 § 7 第 8 号的记号下，

\[
(\mathbf {u} + \mathbf {v}) ^ {\alpha} = \sum_ {\beta + \gamma = \alpha} \frac {\alpha !}{\beta ! \gamma !} \mathbf {u} ^ {\beta} \mathbf {v} ^ {\gamma}.
\]

对于 \(x\in \mathbf{X}\)

\[
(u _ {x} + v _ {x}) ^ {\alpha (x)} = \sum_ {m + n = \alpha (x)} \frac {\alpha (x) !}{m ! n !} u _ {x} ^ {m} v _ {x} ^ {n}
\]

由推论 1 得。将这些方程对 \(x \in \mathbf{X}\) 取乘积，并利用 (13)，即得推论。

命题 2. 设 A 是一个环， \(x_{1}, \ldots, x_{n}\) 是 A 的元素，且 I = \{1, 2, \ldots, n\}。对 H ⊂ I，我们记 \(x_{\mathrm{H}} = \sum_{i \in \mathrm{H}} x_{i}\) 。则

\[
(- 1) ^ {n} \sum_ {\sigma \in \mathfrak {S} _ {n}} x _ {\sigma (1)} \dots x _ {\sigma (n)} = \sum_ {H \subset I} (- 1) ^ {\operatorname{Card} H} (x _ {H}) ^ {n}.\tag{16}
\]

特别地，若 A 是交换的，

\[
(- 1) ^ {n} n! x _ {1} x _ {2} \dots x _ {n} = \sum_ {\mathrm{H} \in \mathrm{I}} (- 1) ^ {\text { Card   H }} (x _ {\mathrm{H}}) ^ {n}.
\]

设 C 为 I 到 \(\{0,1\}\) 中的映射组成的集合。若把每个 \(H \subset I\) 映到它的特征函数，则得到 \(\mathfrak{P}(I)\) 到 C 上的一个双射。因此 (16) 的右边等于：

\[
\begin{array}{r l} \sum_ {a \in \mathbb {C}} (- 1) ^ {a (1) + \dots + a (n)} & \left(\sum_ {i \in \mathbb {I}} a (i) x _ {i}\right) ^ {n} \\ & = \sum_ {a \in \mathbb {C}} (- 1) ^ {a (1) + \dots + a (n)} \sum_ {(i _ {1}, \ldots , i _ {n}) \in \mathbb {I} ^ {n}} a (i _ {1}) \ldots a (i _ {n}) x _ {i _ {1}} \ldots x _ {i _ {n}} \\ & = \sum_ {(i _ {1}, \ldots , i _ {n}) \in \mathbb {I} ^ {n}} c _ {i _ {1} \ldots i _ {n}} x _ {i _ {1}} \ldots x _ {i _ {n}} \end{array}
\]

其中

\[
c _ {i _ {1} \dots i _ {n}} = \sum_ {a \in \mathbb {C}} (- 1) ^ {a (1) + \dots + a (n)} a (i _ {1}) \dots a (i _ {n}).
\]

\begin{enumerate}
\def\labelenumi{(\arabic{enumi})}
\tightlist
\item
  假设 \((i_{1},\ldots,i_{n})\) 不是 I 的一个排列。存在一个 \(j\in I\) ，它不同于 \(i_{1},\ldots,i_{n}\) 。设 \(C'\) 为由 \(a\in C\) （满足 \(a(j)=0\) ）组成的集合。对所有 \(a\in C'\) ，设 \(a^{*}\) 为 a 与 \(\{j\}\) 的特征函数之和。则 \(a^{*}(1)+\cdots+a^{*}(n)=a(1)+\cdots+a(n)+1\) ，因此
\end{enumerate}

\[
\begin{array}{l} c _ {i _ {1} \dots i _ {n}} = \sum_ {a \in C ^ {\prime}} (- 1) ^ {a (1) + \dots + a (n)} a (i _ {1}) \dots a (i _ {n}) + (- 1) ^ {a ^ {*} (1) + \dots + a ^ {*} (n)} a ^ {*} (i _ {1}) \dots a ^ {*} (i _ {n}) \\ = \sum_ {a \in C ^ {\prime}} ((- 1) ^ {a (1) + \dots + (n)} + (- 1) ^ {a (1) + \dots + a (n) + 1}) a (i _ {1}) \dots a (i _ {n}) = 0. \end{array}
\]

\begin{enumerate}
\def\labelenumi{(\arabic{enumi})}
\setcounter{enumi}{1}
\tightlist
\item
  假设存在 \(\sigma \in \mathfrak{S}_n\) ，使得 \(i_1 = \sigma(1), \ldots, i_n = \sigma(n)\) 。则 \(a(i_1) \ldots a(i_n) = 0\) ，除非 \(a\) 仅取值 1。因此 \(c_{i_1 \ldots i_n} = (-1)^n\) 。
\end{enumerate}

